% ============================================================
% CAPÍTULO 8 — Recomendaciones para Futuros Proyectos
% ============================================================
\chapter{Recomendaciones para Futuros Proyectos}
\label{cap:recomendaciones}

A modo de cierre, este cap\'itulo recoge seis recomendaciones derivadas
de la experiencia del Proyecto Demeter, aplicables tanto a futuras
iniciativas de la Asociaci\'on ESIBot como a otros proyectos docentes
con caracter\'isticas similares. Constituyen el \'ultimo activo de
valor del proyecto: la capitalizaci\'on de la experiencia para terceros.

\section{Dimensionar al equipo realmente disponible}

No dise\~nar arquitecturas que requieran m\'as equipo del que se puede
sostener realmente. La ambici\'on t\'ecnica debe escalar con la
capacidad del equipo, no inversamente. Un proyecto con tres componentes
bien hechos por un equipo de tres personas suele aportar m\'as valor
que un proyecto con diez componentes mal mantenidos por una sola.

\section{Distribuir el conocimiento desde el inicio}

\textit{Pair programming}, \textit{code reviews} sistem\'aticas,
documentaci\'on viva y rotaci\'on de responsabilidades. El \textit{bus
factor} debe ser superior a uno desde el primer d\'ia. Si el conocimiento
de un componente del sistema vive en una sola cabeza, ese componente
es fr\'agil estructuralmente, independientemente de su calidad t\'ecnica.

\section{Definir criterios de cierre desde el inicio}

Cada proyecto deber\'ia tener criterios expl\'icitos de:

\begin{itemize}[leftmargin=*]
    \item \textbf{\'Exito}: \textquestiondown cu\'ando se considera completado?
    \item \textbf{Pivote}: \textquestiondown cu\'ando cambiar de
    direcci\'on o redefinir el alcance?
    \item \textbf{Cierre}: \textquestiondown cu\'ando concluir formalmente
    si las condiciones de viabilidad ya no se cumplen?
\end{itemize}

Estos criterios deber\'ian estar documentados antes de empezar y
revisarse peri\'odicamente. Su ausencia es una de las causas principales
de proyectos que se prolongan en estado moribundo.

\section{MVP iterativo $>$ arquitectura aspiracional}

Empezar con la versi\'on m\'as simple que funcione \textit{end-to-end} es
preferible a dise\~nar la arquitectura completa por anticipado. Iterar
a\~nadiendo complejidad solo cuando se justifica con uso real, no con
anticipaci\'on. Esta es una de las lecciones m\'as repetidas en la
literatura de \textit{lean startup} y se aplica igual a proyectos
docentes que a empresariales.

\section{Documentar para audiencias distintas}

La documentaci\'on t\'ecnica para evaluadores no es la misma que la
documentaci\'on interna para futuros mantenedores. La primera enfatiza
arquitectura, decisiones y resultados; la segunda enfatiza c\'omo hacer
cosas concretas, c\'omo levantar el proyecto, c\'omo debuggear. Crear
ambas desde el inicio --- y no como tarea final --- multiplica la
probabilidad de que el proyecto sobreviva la rotaci\'on de miembros.

\section{Reconocer los l\'imites de la asociaci\'on estudiantil}

Algunas iniciativas requieren estructuras m\'as permanentes que una
asociaci\'on universitaria: grupos de investigaci\'on, \textit{spin-offs},
o convenios formales de larga duraci\'on. Una asociaci\'on es excelente
para aprender y prototipar; menos adecuada para operar producci\'on a
largo plazo. Identificar tempranamente cu\'ando un proyecto excede el
formato asociativo permite buscar marcos institucionales m\'as
adecuados antes de generar deuda estructural.

\bigskip

\begin{conclusionbox}
El Proyecto Demeter ha cumplido sus objetivos did\'acticos: ha generado
competencias t\'ecnicas en sus miembros, documentaci\'on reutilizable,
un sistema funcional verificable y --- como cierre --- un conjunto de
lecciones aprendidas que beneficiar\'an a futuros proyectos. Su cierre
formal en este momento, lejos de constituir un fracaso, es la
expresi\'on madura del an\'alisis de viabilidad realizado con la
honestidad t\'ecnica que un Proyecto Docente requiere.
\end{conclusionbox}
